\documentclass[aspectratio=169]{beamer}

% Theme selection - Metropolis is clean and professional for academic presentations
\usetheme{metropolis}
\usecolortheme{default} % Uses red/grey tones, distinct and readable

% Packages
\usepackage{booktabs} % For nicer tables
\usepackage{amsmath}
\usepackage{graphicx}
\usepackage{tikz}

% Title Information
\title[Incentivizing Fairness in BitTorrent]{Incentivizing Fairness in BitTorrent}
\subtitle{An Exposition of Mechanism Design and Tokenomics}
\author{Sravanth Potluri}
\institute{University of Virginia}
\date{December 8th, 2025}

\begin{document}

%--- SLIDE 1: Title ---
\begin{frame}
    \titlepage
\end{frame}

%--- SLIDE 2: Introduction ---
\begin{frame}{The Economic Imperative of P2P}
    \begin{block}{The Problem: Tragedy of the Commons}
        \begin{itemize}
            \item Early P2P (e.g., Gnutella) relied on voluntary altruism.
            \item Rational agents have a dominant strategy to "free-ride": consume without contributing.
            \item Result: By 2005, the vast majority of Gnutella users shared zero files.
        \end{itemize}
    \end{block}

    \vspace{0.5cm}
    \textbf{The BitTorrent Solution (2003):}
    \begin{itemize}
        \item Replaced altruism with \textbf{Direct Reciprocity} (Tit-for-Tat).
        \item Engineered an endogenous incentive mechanism where contribution is the only rational strategy.
    \end{itemize}
\end{frame}

%--- SLIDE 3: The Reference Mechanism ---
\begin{frame}{The Baseline Mechanism: Cohen's Design}
    BitTorrent relies on two core algorithms to manage distribution:

    \begin{columns}
        \column{0.5\textwidth}
        \textbf{1. Piece Selection (Rarest First)}
        \begin{itemize}
            \item Prioritizes pieces least replicated in the swarm.
            \item \textbf{Availability:} Preserves the file if seeders leave.
            \item \textbf{Trade Efficiency:} Ensures peers hold data others want.
        \end{itemize}

        \column{0.5\textwidth}
        \textbf{2. Peer Selection (Choking)}
        \begin{itemize}
            \item Unchokes top $k$ (usually 4) peers providing the highest download rates.
            \item \textbf{Optimistic Unchoking:} Randomly unchokes one neighbor every 30s to discover better partners and bootstrap new peers.
        \end{itemize}
    \end{columns}
\end{frame}

%--- SLIDE 4: Strategic Vulnerabilities ---
\begin{frame}{Strategic Vulnerabilities: BitTyrant}
    Is the reference implementation a Nash Equilibrium? \textbf{No.}
    
    \begin{alertblock}{The "Equal Split" Inefficiency }
        Standard clients split upload bandwidth equally among unchoked peers ($U_{cap}/k$). This leads to "overpayment" or wasted altruism.
    \end{alertblock}

    \textbf{The BitTyrant Strategy:}
    Maximize \textit{Reciprocation Density} by ranking peers via benefit-to-cost ratio:
    \begin{equation}
        Rank_{j} = \frac{d_{j}}{u_{j}}
    \end{equation}
    Where $d_j$ is expected download rate and $u_j$ is minimum upload required to be unchoked.
    
    \begin{itemize}
        \item \textbf{Result:} Median 70\% performance improvement.
        \item \textbf{Risk:} If adopted universally, new peers starve.
    \end{itemize}
\end{frame}

%--- SLIDE 5: Mechanism Design Theory ---
\begin{frame}{BitTorrent as an Auction}
    Levin et al. (2008) reframed the Choking Algorithm as a \textbf{Tournament Auction}.
    
    \begin{itemize}
        \item \textbf{The Flaw:} You only need to beat the $(k+1)^{th}$ bidder to win. Bidding higher yields no extra reward.
    \end{itemize}

    \vspace{0.5cm}
    \textbf{The Theoretical Fix: Proportional Share (PropShare)}
    Allocate bandwidth proportional to contribution:
    \begin{equation}
        U_{i\rightarrow j} = U_{cap} \times \frac{D_{j\rightarrow i}}{\sum_{m\in N} D_{m\rightarrow i}}
    \end{equation}
    
    \begin{itemize}
        \item \textbf{Theorem:} Truth-telling (max upload) becomes a dominant strategy.
        \item \textbf{Reality:} Hard to implement due to TCP fragmentation.
    \end{itemize}
\end{frame}

%--- SLIDE 6: The Seeder's Dilemma ---
\begin{frame}{The Seeder's Dilemma \& Long-Term Incentives}
    \textbf{The Problem:} Once a file is complete, a rational peer has no "want" to trade and should disconnect ("hit-and-run").
    
    \vspace{0.5cm}
    \begin{columns}
        \column{0.48\textwidth}
        \begin{block}{Sociological: Private Trackers}
            \begin{itemize}
                \item Enforce Share Ratio ($R > 1$).
                \item \textbf{Credit Squeeze:} Bandwidth becomes surplus; leechers become scarce.
                \item 3-5x faster than public swarms.
            \end{itemize}
        \end{block}

        \column{0.48\textwidth}
        \begin{block}{Cryptographic: Dandelion}
            \begin{itemize}
                \item Decentralized credit exchange.
                \item \textbf{Cross-Swarm Trading:} Seed a Linux ISO to buy a movie.
                \item Uses "virtual credit" to solve double coincidence of wants.
            \end{itemize}
        \end{block}
    \end{columns}
\end{frame}

%--- SLIDE 7: The Tokenomic Era ---
\begin{frame}{The Tokenomic Era: BitTorrent Token (BTT)}
    Introduced in 2019, moving from Barter to \textbf{Market Exchange}.
    
    \textbf{BitTorrent Speed Protocol:}
    \begin{itemize}
        \item Uses a First-Price Auction with explicit monetary bids.
        \item Bid Equation Strategy:
        \begin{equation}
            Bid_{price} = \frac{\text{Spending Rate} \times \text{Balance}}{\text{Remaining Bytes}}
        \end{equation}
    \end{itemize}

    \textbf{Incentivizing the Long Tail:}
    \begin{itemize}
        \item Rare files $\rightarrow$ Higher Bids $\rightarrow$ Signal to Seeders.
        \item Creates a profit motive to preserve rare content.
    \end{itemize}
\end{frame}

%--- SLIDE 8: Comparative Synthesis ---
\begin{frame}{Comparative Synthesis}
    \begin{table}[]
        \resizebox{\textwidth}{!}{%
        \begin{tabular}{@{}lllll@{}}
        \toprule
        \textbf{Feature} & \textbf{Tit-for-Tat (Cohen)} & \textbf{PropShare (Levin)} & \textbf{Dandelion} & \textbf{BTT (Token)} \\ \midrule
        \textbf{Incentive Basis} & Immediate Barter & Proportional Barter & Virtual Credit & Crypto-Currency \\
        \textbf{Seeding Incentive} & Zero (Altruism) & Zero (Altruism) & Future Credit & Profit (BTT) \\
        \textbf{Strategy Proof?} & No (BitTyrant) & Yes (Ideally) & Yes & Market Dependent \\
        \textbf{Rare Content} & Poor Availability & Poor Availability & Good Availability & High Availability \\
        \textbf{Sybil Resistance} & Moderate (IP) & High (Dilution) & High (Costly ID) & High (Costly Token) \\ \bottomrule
        \end{tabular}%
        }
        \caption{Comparison of Incentive Mechanisms }
    \end{table}
\end{frame}

%--- SLIDE 9: Future Directions ---
\begin{frame}{Future Perspectives}
    The strict binary between "Barter" and "Money" is evolving.
    
    \begin{itemize}
        \item \textbf{Reputation DAOs:}
        \begin{itemize}
            \item Non-transferable reputation scores secured by DAGs.
            \item Prevents "pay-to-win" scenarios inherent in fiat-based tokens.
        \end{itemize}
        
        \vspace{0.3cm}
        \item \textbf{Smart Seeders (AI Agents):}
        \begin{itemize}
            \item Autonomous agents that optimize file portfolios.
            \item Crawl DHT for high yield (low supply/high demand) swarms.
            \item Turns file storage into a high-frequency trading market.
        \end{itemize}
    \end{itemize}
\end{frame}

%--- SLIDE 10: Conclusion ---
\begin{frame}{Conclusion}
    \begin{itemize}
        \item \textbf{Evolution:} We have moved from "Communal Sharing" (TFT) to "Market Exchange" (BTT).
        \item \textbf{Lessons Learned:}
        \begin{enumerate}
            \item Immediate barter (TFT) is robust but insufficient for rare content.
            \item Long-term memory (Private Trackers/Credits) is superior to short-term optimization.
            \item Financialization (BTT) introduces friction but solves the Seeder's Dilemma.
        \end{enumerate}
        \item \textbf{The Future:} The intersection of Mechanism Design and Cryptographic Tokenomics.
    \end{itemize}
    
    \vspace{1cm}
    \centering
    \Large \textbf{Thank You. Questions?}
\end{frame}

%--- SLIDE 11: References ---
\begin{frame}[allowframebreaks]{References}
    \bibliographystyle{ieeetr}
    \tiny
    \begin{itemize}
        % The curly braces {} prevent Beamer from misinterpreting the brackets
        \item {[1]} B. Cohen, "Incentives build robustness in BitTorrent," 2003.
        \item {[2]} M. Piatek et al., "Do incentives build robustness in BitTorrent?" NSDI, 2007.
        \item {[3]} D. Levin et al., "BitTorrent is an auction," SIGCOMM, 2008.
        \item {[4]} M. Sirivianos et al., "Dandelion: Cooperative content distribution," 2007.
        \item {[5]} M. Meulpolder et al., "Public and private BitTorrent communities," 2010.
        \item {[6]} BitTorrent Foundation, "BitTorrent (BTT) White Paper," 2019.
    \end{itemize}
\end{frame}

\end{document}